\section{Preliminares matemáticos}

\subsection{Sistemas dinámicos y balances conservativos}

\begin{definition}[Sistema dinámico en tiempo continuo]\label{def11}
Sea $\mathbf{x}(t) \in \mathbb{R}^n$ el vector de \emph{variables de estado},
$\mathbf{u}(t) \in \mathbb{R}^m$ el vector de \emph{entradas} (forzamientos externos)
y $\mathbf{y}(t) \in \mathbb{R}^p$ el vector de \emph{mediciones}. Un sistema dinámico
es el par de ecuaciones
\begin{equation}
\dot{\mathbf{x}}(t) = f\bigl(\mathbf{x}(t), \mathbf{u}(t); \boldsymbol{\theta}\bigr),
\qquad
\mathbf{y}(t) = g\bigl(\mathbf{x}(t); \boldsymbol{\theta}\bigr),
\label{eq:sistema-general}
\end{equation}
donde $f$ describe la dinámica interna, $g$ el mapa de observación y
$\boldsymbol{\theta} \in \mathbb{R}^q$ el vector de parámetros.
\end{definition}

\begin{definition}[Balance conservativo]\label{def12}
Sea $x(t)$ una variable de estado que representa una cantidad extensiva
(masa, energía, moles) contenida en un volumen de control. Su dinámica
obedece el balance
\begin{equation}
\dot{x}(t) = \underbrace{J_{\mathrm{in}}(t)}_{\text{entradas}}
- \underbrace{J_{\mathrm{out}}(t)}_{\text{salidas}}
+ \underbrace{P(t)}_{\text{producción}}
- \underbrace{C(t)}_{\text{consumo}},
\label{eq:balance}
\end{equation}
donde los flujos $J$ cruzan la frontera del volumen de control, mientras que
$P$ y $C$ representan transformaciones \emph{internas} de una forma de materia
o energía a otra.
\end{definition}

\begin{figure}[htbp]
\centering
\begin{tikzpicture}[
  font=\small,
  vol/.style={draw, thick, rounded corners=8pt, minimum width=5.2cm,
              minimum height=3.6cm, fill=blue!6},
  flujo/.style={-{Stealth[length=3mm]}, very thick},
  trans/.style={-{Stealth[length=2.5mm]}, thick}
]

% Volumen de control
\node[vol] (vc) {};
\node at ([yshift=1.15cm]vc.center) {\textbf{Volumen de control}};
\node at ([yshift=0.35cm]vc.center) {$x(t)$};
\node at ([yshift=-0.15cm]vc.center) {\footnotesize (masa o energía acumulada)};

% Frontera
\draw[dashed, thick] ([xshift=6pt,yshift=6pt]vc.north west)
  rectangle ([xshift=-6pt,yshift=-6pt]vc.south east);

% Transformaciones internas: ciclo P y C
\draw[trans, orange!85!black] ([xshift=-1.0cm,yshift=-0.85cm]vc.center)
  to[bend left=25] node[above, pos=0.5]{$P$}
  ([xshift=1.0cm,yshift=-0.85cm]vc.center);
\draw[trans, orange!85!black] ([xshift=1.0cm,yshift=-1.25cm]vc.center)
  to[bend left=25] node[below, pos=0.5]{$C$}
  ([xshift=-1.0cm,yshift=-1.25cm]vc.center);

% Flujos que cruzan la frontera
\draw[flujo, green!50!black] (-5.4,0.35) -- node[above, midway]{$J_{\mathrm{in}}$}
  ([yshift=0.35cm]vc.west);
\draw[flujo, red!70!black] ([yshift=0.35cm]vc.east) -- node[above, midway]{$J_{\mathrm{out}}$}
  (5.4,0.35);

\end{tikzpicture}
\caption{Esquema del balance conservativo de la Definición~1.2. La cantidad
acumulada $x(t)$ dentro del volumen de control cambia por dos mecanismos:
flujos que \emph{cruzan la frontera} ($J_{\mathrm{in}}$, $J_{\mathrm{out}}$)
y transformaciones \emph{internas} ($P$: producción, $C$: consumo) que
convierten una forma de materia o energía en otra sin cruzar la frontera.}
\label{fig:balance}
\end{figure}

\begin{remark}
La ecuación \eqref{eq:balance} es la estructura común de todos los modelos de
este documento. El modelo de Eriksen \cite{eriksenDynamicModellingFeed2022} es una red de balances
de masa entre compartimentos ($A$, $B$, $L$) sin términos $J$ (sistema cerrado
a nivel de larva); el modelo del contenedor que desarrollaremos añade los
términos de entrada y salida asociados al flujo de aire (\emph{inlet/outlet})
y a la alimentación.
\end{remark}

\begin{definition}[Compartimento]\label{def13}
Un \emph{compartimento} es una variable de estado $x_i(t)$ que representa la
cantidad total de una entidad homogénea (masa estructural, reserva lipídica,
gas en un volumen de aire) dentro del volumen de control, bajo el supuesto de
\emph{mezcla homogénea}: la entidad está uniformemente distribuida, de modo
que un único valor escalar describe su estado en todo el volumen.
\end{definition}

\begin{figure}[htbp]
\centering
\begin{tikzpicture}[
  font=\small,
  vol/.style={draw, thick, rounded corners=8pt, minimum width=4.4cm,
              minimum height=3.4cm, fill=blue!6},
  sensor/.style={draw, thick, fill=red!15, rounded corners=2pt}
]

% Volumen de control con mezcla homogénea
\node[vol] (vc) {};
\node at ([yshift=1.25cm]vc.center) {\textbf{Compartimento}};

% Partículas uniformemente distribuidas (mezcla homogénea)
\foreach \p in {(-1.5,0.6),(-0.9,0.1),(-0.3,0.7),(0.3,0.2),(0.9,0.8),(1.5,0.3),
                (-1.4,-0.5),(-0.7,-0.9),(0.0,-0.4),(0.7,-1.0),(1.4,-0.6),(-1.0,-1.2),(0.4,-0.9)}
  \fill[blue!55!black] \p circle (1.6pt);

% Estado escalar
\node at ([yshift=-1.5cm]vc.center) {$x_i(t)$};

% Sensor puntual
\node[sensor, minimum width=0.5cm, minimum height=0.35cm] (s) at (vc.east) {};
\node[right=0.15cm of s] {$y_i(t)$};
\draw[-{Stealth[length=2mm]}, thick, red!70!black] (s.west) -- ++(-0.45,0);

% Anotación del supuesto
\node[below=0.35cm of vc, text width=5.2cm, align=center]
  {\footnotesize Supuesto de mezcla homogénea: cualquier punto del volumen
   representa al todo.};

\end{tikzpicture}
\caption{Interpretación del supuesto de mezcla homogénea (Definición~1.3).
La entidad acumulada se supone uniformemente distribuida en el volumen, de
modo que el estado escalar $x_i(t)$ la describe por completo y una medición
puntual $y_i(t)$ (por ejemplo, un sensor de CO$_2$ o de temperatura) es
representativa de todo el compartimento.}
\label{fig:compartimento}
\end{figure}

La Figura~\ref{fig:compartimento} ilustra el contenido del supuesto de mezcla
homogénea. Los puntos distribuidos uniformemente en el volumen representan la
entidad acumulada ---gas, biomasa o humedad--- y sintetizan el supuesto sin
necesidad de ecuaciones: al no existir gradientes internos, un único escalar
$x_i(t)$ describe el estado de todo el compartimento. El elemento de medición
puntual $y_i(t)$ representa la conexión con el sistema experimental: la
lectura de un sensor (de temperatura, humedad relativa o concentración de
CO$_2$) ubicado en una posición fija del contenedor se toma como
representativa del volumen completo. Este es precisamente el supuesto que
deberá vigilarse con mayor cuidado en la etapa de validación, pues en el
sistema real cabe esperar gradientes locales de temperatura y gases en la
vecindad de la masa larval, donde la actividad metabólica es intensa.

\begin{definition}[Red de compartimentos]\label{def14}
Una \emph{red de compartimentos} es un conjunto $\{x_1, \dots, x_n\}$ acoplado
por flujos de transferencia $J_{ij}(t) \geq 0$ del compartimento $j$ al
compartimento $i$, de modo que
\begin{equation}
\dot{x}_i = \sum_{j \neq i} \bigl( J_{ij} - J_{ji} \bigr)
+ J_{i,\mathrm{in}} - J_{i,\mathrm{out}} + P_i - C_i,
\qquad i = 1, \dots, n.
\label{eq:red}
\end{equation}
Las transferencias internas satisfacen conservación global: si $P_i = C_i = 0$
y no hay flujos externos, entonces $\sum_i \dot{x}_i = 0$.
\end{definition}

\begin{remark}
El modelo de Eriksen es una red de tres compartimentos ---asimilados $A$,
biomasa estructural $B$ y lípidos $L$--- con transferencias internas
$J_{AB}, J_{AL}, J_{LB}$ (asimilación, síntesis y reciclaje). El sistema
completo alimento--larvas--ambiente extiende esta red con compartimentos
externos: masa de sustrato, y moles de CO$_2$, CH$_4$ y vapor de agua en el
volumen de aire del contenedor.
\end{remark}
\begin{figure}[htbp]
\centering
\begin{tikzpicture}[
  font=\small,
  comp/.style={draw, thick, circle, minimum size=1.15cm, fill=blue!8},
  ext/.style={draw, thick, rounded corners=3pt, fill=green!10,
              minimum width=1.5cm, minimum height=0.85cm},
  flujo/.style={-{Stealth[length=2.6mm]}, thick},
  interno/.style={-{Stealth[length=2.6mm]}, thick, orange!85!black},
  resp/.style={-{Stealth[length=2.6mm]}, thick, red!70!black}
]

% Compartimentos internos (modelo de Eriksen)
\node[comp] (A) at (0,2.6) {$A$};
\node[comp] (B) at (-1.9,0) {$B$};
\node[comp] (L) at (1.9,0) {$L$};

% Nodo externo: sustrato (forzamiento)
\node[ext] (S) at (0,5.0) {Sustrato};

% Flujo externo de entrada
\draw[flujo, green!50!black] (S) -- node[right, pos=0.45]{$r_A = aB$
  \footnotesize (asimilación)} (A);

% Transferencias internas
\draw[interno] (A) to[bend right=18] node[left, pos=0.45]{$r_B$} (B);
\draw[interno] (A) to[bend left=18] node[right, pos=0.45]{$r_L$} (L);
\draw[interno] (L) to[bend left=15] node[right, pos=0.4]
  {\footnotesize reciclaje ($r_L < 0$)} (A);

% Tres fuentes de CO2
\draw[resp] (B.south) -- ++(0,-1.0)
  node[below]{$r_{\mathrm{CO_2},m} = mB$};
\draw[resp] (-1.35,1.35) -- ++(-0.85,0.85)
  node[above left=-2pt]{$r_{\mathrm{CO_2},B} = Y_B r_B$};
\draw[resp] (1.35,1.35) -- ++(0.85,0.85)
  node[above right=-2pt]{$r_{\mathrm{CO_2},L} = Y_L r_L$};

% Etiquetas de los compartimentos
\node[left=0.35cm of A, align=right] {\footnotesize asimilados};
\node[left=0.35cm of B, align=right] {\footnotesize biomasa\\[-1pt]\footnotesize estructural};
\node[right=0.35cm of L, align=left] {\footnotesize lípidos\\[-1pt]\footnotesize de reserva};

\end{tikzpicture}
\caption{El modelo de Eriksen \cite{eriksenDynamicModellingFeed2022} como red de compartimentos
(Definición~1.3). El sustrato es un forzamiento externo que alimenta el
compartimento de asimilados $A$ a tasa $r_A = aB$; las flechas internas
representan síntesis de biomasa estructural ($r_B$), acumulación de lípidos
($r_L$) y reciclaje de reservas hacia el pool de asimilados cuando
$r_L < 0$. Las tres flechas rojas son las fuentes de producción de CO$_2$:
mantenimiento (proporcional a $B$), costo de síntesis estructural y costo de
síntesis de lípidos, cuya suma es la producción total
(Ec.~\eqref{eq:rco2total}).}
\label{fig:red-eriksen}
\end{figure}

En el modelo de Eriksen, la producción total de CO$_2$ se descompone en tres
procesos metabólicos distintos \cite{eriksenDynamicModellingFeed2022}. El primero es el
\emph{mantenimiento}: la energía que la larva requiere para sostener los
procesos vitales que no producen crecimiento neto, modelada como proporcional
a la biomasa estructural,
\begin{equation}
r_{\mathrm{CO_2},m} = m\,B,
\label{eq:rco2m}
\end{equation}
donde $m$ es la tasa específica de mantenimiento (en el artículo,
$m = 0.08$~día$^{-1}$). El segundo es el \emph{costo de síntesis de la
biomasa estructural}: construir tejido a partir de asimilados es un proceso
energéticamente costoso, y esa energía se obtiene respirando asimilados,
\begin{equation}
r_{\mathrm{CO_2},B} = Y_B\,r_B,
\label{eq:rco2b}
\end{equation}
con $Y_B = 0.44$ el costo (adimensional) de crecimiento estructural y $r_B$
la tasa de síntesis de biomasa. El tercero es el \emph{costo de síntesis de
lípidos de reserva},
\begin{equation}
r_{\mathrm{CO_2},L} = Y_L\,r_L,
\label{eq:rco2l}
\end{equation}
con $Y_L = 0.42$ y $r_L$ la tasa de acumulación de lípidos; este término se
anula cuando la asimilación no cubre el mantenimiento y los lípidos se
reciclan en lugar de sintetizarse ($r_L < 0$). La producción total es la suma
de los tres procesos:
\begin{equation}
r_{\mathrm{CO_2}} = r_{\mathrm{CO_2},m} + r_{\mathrm{CO_2},B} + r_{\mathrm{CO_2},L}.
\label{eq:rco2total}
\end{equation}
En la Figura~\ref{fig:red-eriksen}, las flechas de salida de CO$_2$
corresponden a estos tres términos. Esta descomposición es la conexión directa
con el sistema instrumentado: la señal de CO$_2$ medida en el \emph{outlet}
del contenedor contiene, superpuestas, las tres fuentes, de modo que un modelo
como el anterior permite inferir qué fracción de la señal corresponde a
mantenimiento, a crecimiento estructural o a acumulación de lípidos.

\subsection{El modelo de Eriksen como sistema dinámico}

El modelo de Eriksen \cite{eriksenDynamicModellingFeed2022} se obtiene al particularizar la
Definición~1.1 con el vector de estado, las entradas y las mediciones
\begin{equation}
\mathbf{x} = \begin{bmatrix} A \\ B \\ L \end{bmatrix},
\qquad
\mathbf{u} = \begin{bmatrix} s \\ T \end{bmatrix},
\qquad
\mathbf{y} = r_{\mathrm{CO_2}},
\label{eq:eriksen-sistema}
\end{equation}
donde $s(t)$ representa el suministro de alimento y $T(t)$ la temperatura de
cría. La función de dinámica $f$ está dada por los balances de cada
compartimento,
\begin{equation}
\dot{A} = r_A - r_B - r_L - r_{\mathrm{CO_2}},
\qquad
\dot{B} = r_B,
\qquad
\dot{L} = r_L,
\label{eq:eriksen-edos}
\end{equation}
donde las tasas $r_A$, $r_B$, $r_L$ y $r_{\mathrm{CO_2}}$ son funciones del
estado y de los parámetros
$\boldsymbol{\theta} = (a, m, Y_B, Y_L, B_{\max}, \beta, \dots)$.
El mapa de observación $g$ es la Ec.~\eqref{eq:rco2total}: la producción
total de CO$_2$ es una medición indirecta de la actividad metabólica de los
tres compartimentos, ninguno de los cuales es medible directamente de forma
no destructiva.

\begin{remark}[Cierre: del sistema general al modelo de Eriksen]
La formulación \eqref{eq:eriksen-sistema}--\eqref{eq:eriksen-edos} pone de
manifiesto dos extensiones que motivan el presente documento. Primero, en el
modelo original la temperatura $T$ y el suministro $s$ son forzamientos
exógenos; en un contenedor cerrado e instrumentado, la temperatura es un
estado más del sistema, acoplado al calor metabólico que las propias larvas
generan. Segundo, el mapa de observación $g$ puede enriquecerse con las
mediciones de humedad relativa y CH$_4$, lo que plantea el problema de
estimación de estados y parámetros que abordaremos con redes neuronales
informadas por física. En síntesis, esta sección ha establecido el lenguaje
formal sobre el que se construye el resto del documento: un sistema dinámico
(Definición~\ref{def11}) cuyas ecuaciones son balances conservativos
(Definición~\ref{def12}) organizados en una red de compartimentos
(Definiciones~\ref{def13} y~\ref{def14}), bajo el supuesto de mezcla
homogénea ilustrado en la Figura~\ref{fig:compartimento}. El modelo de
Eriksen \cite{eriksenDynamicModellingFeed2022} emerge así como un caso particular de esta
estructura ---una red de tres compartimentos con tres fuentes de producción
de CO$_2$--- y su reformulación en términos de estados, entradas y
mediciones delimita con precisión el punto de partida de la extensión que se
desarrolla en las secciones siguientes: acoplar la dinámica larval con la
del ambiente del contenedor y con un mapa de observación enriquecido por
instrumentación en tiempo real.
\end{remark}

